\documentclass{situationist-preprint}

\title{Post-Mortem Survival Under a Pattern Theory of Persons}
\msid{REL-08-a}
\msfamily{REL-08}
\msdiscipline{Philosophy of Religion}
\mstype{Principal manuscript}
\mskeywords{life after death; resurrection; pattern theory; psychological continuity; replica theory; duplication objection; animalism; mediumship; evidential inaccessibility}

\begin{document}
\maketitle

\begin{abstract}
Is there life after death? We separate two questions: whether survival of death is coherent, and whether any evidence available before death could confirm it. On the first, we argue that survival is coherent under a pattern theory of persons, on which a person is a concrete realization of a psychological pattern and persists through the reliably caused, non-branching continuation of that pattern, whatever its substrate and across temporal gaps. We give a model of post-mortem survival satisfying the theory and show that the standard objections, van Inwagen's argument from the burned manuscript and the duplication objection to Hick's replica theory, presuppose persistence conditions the theory rejects. Every case of survival under the theory meets the temporal condition that governs teletransportation. On the second question, we prove that the facts which, on the theory, distinguish survival from the mere post-mortem realization of one's pattern, namely the absence of a second realization and the right kind of causal link, make no difference to anything observable before death. Evidence may bear on whether one's pattern is realized after death; it cannot favor survival over its matched rivals. We close by relating the result to the animalist conclusion of PID-05.
\end{abstract}

\section{The problem}

The question whether there is life after death contains two questions that the tradition has not always kept apart. The first is a question of coherence: is there a consistent account of what a person is on which a person who dies could exist, and be the same person, at a later time? The second is a question of evidence: if the answer to the first is yes, is there, or could there be, evidence available to the living that one does survive? We argue that the first question has a positive answer under a pattern theory of persons, and that the second has a negative one. The two answers are connected: the features that make survival coherent on the pattern theory are the features that place it beyond the reach of evidence.

The positions in the literature can be ordered by what they take to survive. On the dualist view of Plato's \emph{Phaedo} and of Descartes, a person is, or has as her essential part, an immaterial soul that is not destroyed by the death of the body; survival is the soul's continued existence \citep{Swinburne1997}. On the hylomorphic view of Aquinas, the soul survives death as the form of the body, but it is not the person: Aquinas says in his commentary on the first letter to the Corinthians that the soul, being a part of the human being, is not the whole human being, and that my soul is not I \citep[15, lect.~2]{AquinasCor}. The person exists again only at the resurrection, when the soul informs a body once more. Materialists who accept a resurrection must explain how a material person destroyed at death can exist later. \citet{vanInwagen1978} argued that reassembly of the same matter would not suffice, and proposed, as a possibility rather than a doctrine, that God might remove the body at death and replace it with a simulacrum. \citet{Zimmerman1999} proposed that at death the elements of the body might undergo fission, one set of successors forming the corpse and the other the resurrected body, with only the latter standing in the immanent causal relation that persistence requires. \citet{Baker2000} and \citet{Corcoran2006} defended constitution views on which a person is constituted by, but not identical to, an organism, and can be constituted later by a different body if it bears the same first-person perspective.

The view we examine descends from Locke and from \citet{Hick1976}. Locke held that personal identity consists in sameness of consciousness, and that this explains how the same person could be raised in a different body \citep[II.xxvii]{Locke1694}. Hick argued that if a person were to cease to exist in one place and an exact psychophysical replica, with the same memories and character, were to appear elsewhere, we would rightly count the replica as the same person, and that the resurrection can be understood as the divine creation of such a replica in a different space. \citet[ch.~6]{Polkinghorne2002} proposed that the soul is the information-bearing pattern carried by the body, and that God, remembering the pattern, can re-embody it. \citet{Steinhart2014} developed a computational version of the same idea. We call the general position a \emph{pattern theory}: persons are realizations of psychological patterns, and survival is the continuation of a pattern in an appropriate way.

On the evidential side, the Society for Psychical Research, founded in 1882, collected reports of apparitions and of mediums who seemed to convey information known only to the dead; \citet{James1898} argued that the dependence of consciousness on the brain might be transmissive rather than productive, leaving survival open. \citet{Broad1925} proposed that a \enquote{psychic factor} might survive death without being a person, and combine with the organism of a medium. \citet{Stevenson1974} documented children who reported memories corresponding to the lives of deceased strangers, and prospective studies of cardiac arrest survivors \citep{vanLommel2001} found that a minority report experiences during the period of arrest. The most careful recent assessments, by \citet{Almeder1992}, who argued that the evidence supports survival, and \citet{Braude2003}, who concluded more guardedly that it favors some form of survival over the hypothesis that the information comes from the living, both turn on the comparison of survival with rivals that explain the same data. We argue that the pattern theory supplies rivals of a kind that these discussions did not consider, and that no evidence can tell against them.

The plan is as follows. Section~\ref{s:history} reviews the debate over the coherence of resurrection. Section~\ref{s:frame} states the pattern theory. Section~\ref{s:coh} proves that survival of death is coherent under it, and Section~\ref{s:cases} applies the theory to the main models of resurrection and to some harder cases. Section~\ref{s:pid} relates the result to the conclusions of the PID family. Section~\ref{s:evid} proves the evidential result, and Section~\ref{s:evcases} applies it to mediumship, reincarnation reports and near-death experiences. Objections, limits and conclusion follow.

\section{The coherence debate}\label{s:history}

The modern debate begins with Locke's claim that the identity of a person consists in consciousness, so that whatever has the consciousness of present and past actions is the same person to whom they belong \citep[II.xxvii.10]{Locke1694}. Locke drew the consequence for the resurrection himself: the same person will be raised, whatever body it has. \citet{Butler1736} objected that consciousness of past actions presupposes personal identity and cannot constitute it, and Reid added the case of the brave officer, who remembers his deeds as a young officer but not the flogging he received as a boy, to show that memory is not transitive while identity is. Both objections were answered in the twentieth century by replacing memory with overlapping chains of psychological connection, caused in the right way \citep{Shoemaker1984,Parfit1984}. On this reformulation the brave officer is the boy, since he is linked to him by a chain, and the circularity objection is met by defining the relevant memories as quasi-memories, states caused in the right way by earlier experiences without presupposing that the subject is the same.

The resurrection debate proper turned on three problems. The first is the \emph{gap problem}. If a person ceases to exist at death and exists again at the resurrection, she has a gappy existence, and some have held that nothing can begin to exist twice. \citet{vanInwagen1978} gave the problem its sharpest form. If a manuscript written by Augustine is burned, God could create a manuscript exactly like it, even of the same atoms reassembled; but the result would not be Augustine's manuscript, only a perfect copy, since it would not stand in the right causal relation to the original. Van Inwagen concluded that a materialist who believes in the resurrection must hold that the body is not in fact destroyed at death, and offered the simulacrum story as a way it might not be.

The second is the \emph{duplication problem}. If God can create one replica of a dead person, God can create two. Both would stand in exactly the relation to the dead person that a single replica would. They cannot both be the dead person, since they are not each other; and since each stands in the same relation to her as the other, neither is. But a single replica stands in that same relation; so, the objection concludes, a single replica would not be the dead person either. \citet{Penelhum1970} pressed this against replica theories, and the argument relies on the principle, which \citet{Wiggins1980} and others have defended, that whether $x$ is identical to $y$ can depend only on $x$, $y$ and the relations between them.

The third is the \emph{causal problem}. If identity over time requires an immanent causal connection between earlier and later stages, a connection running through the states of the persisting thing itself, then a replica created by God from God's knowledge of the dead person, with no causal link to her remains, is not connected to her in the right way. \citet{Zimmerman1999} accepted the requirement and built his falling elevator model to meet it; \citet{Hudson2001} examined related proposals within a four-dimensionalist framework.

Hick's own reply to the duplication problem was that a replica would be the person if and only if there were only one; he accepted that identity would then depend on whether a second replica existed \citep{Hick1976}. Parfit reached the same structure from the other direction: the non-branching psychological criterion makes identity depend on the absence of a competitor, and Parfit took this as a reason to think that identity is not what matters \citep[ch.~12]{Parfit1984}. Constitution views avoid the duplication problem by making personal identity primitive, at the cost of leaving it unexplained what makes a later body constitute the same first-person perspective \citep{Baker2000}. The pattern theory, as we formulate it, accepts Hick's and Parfit's reply and drops the immanent-causation requirement in favor of a requirement of reliable causation. We show that, so formulated, it answers all three problems.

\section{The pattern theory}\label{s:frame}

\begin{definition}[Pattern]\label{d:pattern}
A \emph{psychological pattern} is a structure of psychological states and dispositions: memories and quasi-memories, beliefs, values, intentions, skills and traits of character, together with the functional relations among them by which they jointly produce thought and action.
\end{definition}

\begin{definition}[Realization]\label{d:real}
A concrete system \emph{realizes} a pattern at a time if at that time it is in states that play the functional roles the pattern specifies, so that it thinks and acts as the pattern requires.
\end{definition}

A pattern is a type; a realization is a concrete system at a time. The distinction matters, since persons are not types. On the account of MND-08, thinking is a functional property that systems of different constitution can have, and Definition~\ref{d:real} accordingly does not restrict the substrate of a realization. Nor does it require that a pattern exist unrealized between realizations; a pattern is a structure of a realizer, and on the conclusion of MET-26 information requires a physical medium and structure presupposes something structured. What the theory requires between realizations is not that the pattern float free, but that there be a causal chain.

\begin{definition}[Pattern continuity]\label{d:cont}
A realization $y$ at $t_2$ is \emph{pattern-continuous} with a realization $x$ at $t_1$ if the pattern $y$ realizes is linked to the pattern $x$ realizes by overlapping chains of strong psychological connection, and the states of $y$ depend on those of $x$ through a reliable causal process, that is, a process that would have transmitted a different pattern had $x$ realized a different one.
\end{definition}

The reliability clause is the wide version of Parfit's causal condition, strengthened by a requirement of counterfactual sensitivity. It excludes deviant cases, such as a replica produced by chance or by a designer who did not take the original as a model, while admitting causal chains that do not run through a single body, such as the transmission of a scan.

\begin{principle}[Persons]\label{p:persons}
A person is a concrete realization of a psychological pattern rich enough to include self-consciousness and the capacity to think of oneself as persisting.
\end{principle}

\begin{principle}[Persistence]\label{p:persist}
A person $x$ existing at $t_1$ is identical to a person $y$ existing at $t_2$ if and only if $y$ at $t_2$ is pattern-continuous with $x$ at $t_1$ and the continuity does not branch: there are no two realizations, existing at the same time or at different times, each pattern-continuous with $x$ at $t_1$, neither of which is pattern-continuous with the other.
\end{principle}

Principles~\ref{p:persons} and~\ref{p:persist} together are what we call \emph{the pattern theory}. Principle~\ref{p:persist} is the non-branching criterion of \citet{Shoemaker1984} and \citet{Parfit1984}, adopted in PID-02-a, with the wide causal condition of Definition~\ref{d:cont}. It does not require that the person exist at every time between $t_1$ and $t_2$. On the four-dimensionalist account of PID-11, a person is a sum of temporal parts unified by a persistence relation, and Principle~\ref{p:persist} says which relation; a sum of parts unified by it may be scattered in time as a sum of parts unified by spatial relations may be scattered in space.

\begin{definition}[Death; survival]\label{d:death}
A person \emph{dies} at $t_d$ if the organism that realizes her pattern ceases irreversibly at $t_d$ to realize it. She \emph{survives death} if there is a person existing at some $t > t_d$ who is identical to her.
\end{definition}

\begin{definition}[Coherence]\label{d:coh}
Survival of death is \emph{coherent under} a theory of persons if the theory is consistent with the existence of a person who dies and survives death, and the case violates no principle of persistence that the theory itself accepts.
\end{definition}

Coherence under a theory is weaker than possibility simpliciter. If the pattern theory were necessarily false, survival of death could be coherent under it and yet impossible for us. We return to this in Section~\ref{s:pid}.

\section{Coherence}\label{s:coh}

\begin{lemma}[Gaps]\label{l:gap}
Principle~\ref{p:persist} is consistent with a person's existing at $t_1$ and at $t_2$ and at no time between.
\end{lemma}

\begin{proof}
Pattern continuity, by Definition~\ref{d:cont}, requires psychological connection between the patterns realized at $t_1$ and $t_2$ and a reliable causal process linking them. Neither requires a realization at intermediate times. A process that records the state of a realizer and later produces a realizer in the recorded state is a causal process, and if the record would have been different had the state been different, it is reliable in the sense of Definition~\ref{d:cont}. A case in which such a process links $x$ at $t_1$ to $y$ at $t_2$, with no realization of the pattern between, and with no second realization at any time, satisfies the right-hand side of Principle~\ref{p:persist}.
\end{proof}

Lemma~\ref{l:gap} is not peculiar to post-mortem cases. In the Simple Teletransportation case of PID-02-a, the person is destroyed by the scan and her replica activated after transmission; on the verdict of PID-02-a she survives, and between scan and activation she is not realized anywhere. Survival of death, on the pattern theory, is teletransportation across a longer interval and through a different channel.

\begin{lemma}[Bridge]\label{l:bridge}
Under the pattern theory, the causal process required by Definition~\ref{d:cont} may run through the knowledge and will of an agent.
\end{lemma}

\begin{proof}
Definition~\ref{d:cont} requires that the later states depend on the earlier ones through a process that would have transmitted a different pattern had the earlier realizer realized a different one. Suppose an agent knows the pattern a person realizes at the time of her death, and later produces a realizer of that pattern because she knows it. Had the person realized a different pattern, the agent, knowing that pattern, would have produced a realizer of it. The process is therefore counterfactually sensitive in the required way. Nothing in Definition~\ref{d:cont} requires that the process be immanent in the person or her remains.
\end{proof}

The process of Lemma~\ref{l:bridge} need not be one in which the pattern exists as free-floating information. In the theistic case, the agent is God, whose knowledge, if God is timeless as REL-19 argues, is not located in the interval between death and resurrection at all; the dependence of God's knowledge on the person's pattern is the counterfactual dependence that REL-05 finds between God's knowledge and free choices. What the bridge requires is dependence, not storage.

\begin{theorem}[Coherence]\label{t:coh}
Survival of death is coherent under the pattern theory.
\end{theorem}

\begin{proof}
We describe a case and verify it against each clause of the theory. A person $x$ is realized by an organism $O$ until $t_d$, when $O$ ceases irreversibly to realize her pattern; so $x$ dies at $t_d$ by Definition~\ref{d:death}. An agent who knows the pattern $O$ realized at $t_d$ produces, at $t_r > t_d$, a body $B$ that realizes that pattern, because the agent knows it; no other realization pattern-continuous with $x$ is produced at any time. $B$ at $t_r$ realizes a pattern rich enough for self-consciousness, so by Principle~\ref{p:persons} there is a person $y$ at $t_r$. The pattern $B$ realizes is connected to $x$'s pattern at $t_d$ by the strongest psychological connections, identity of memories, values, intentions and character; the causal link is reliable by Lemma~\ref{l:bridge}; and the gap between $t_d$ and $t_r$ is no obstacle by Lemma~\ref{l:gap}. So $y$ is pattern-continuous with $x$. The continuity does not branch: before $t_d$ only $O$ realizes the pattern, between $t_d$ and $t_r$ nothing does, after $t_r$ only $B$ does, and $B$ is pattern-continuous with $O$'s last realization of it. By Principle~\ref{p:persist}, $y = x$, and by Definition~\ref{d:death}, $x$ survives death. The case uses no principle the theory rejects, so by Definition~\ref{d:coh} survival is coherent under it.
\end{proof}

\begin{proposition}[Timing]\label{p:timing}
Under the pattern theory, a person survives death as a post-mortem realizer $B$ only if $B$ is not activated before the ante-mortem realizer ceases to realize her pattern, and no second realizer pattern-continuous with her is activated at any time.
\end{proposition}

\begin{proof}
If $B$ were activated while the organism still realized the pattern, $B$ and the organism after $B$'s activation would each be pattern-continuous with the person as she was before, and neither would be pattern-continuous with the other, since neither's states depend on the other's; the continuity would branch, and Principle~\ref{p:persist} would be violated. If a second realizer $B'$ were produced from the same source, whether while $B$ existed or after $B$ had ceased to exist, $B$ and $B'$ would each be pattern-continuous with the person at death and neither with the other, and the continuity would again branch.
\end{proof}

Proposition~\ref{p:timing} is the post-mortem form of Theorem~3.3 of PID-02-a: survival by teletransportation requires that the original be destroyed before the replica is activated. Death is the destruction of the original, and survival of death satisfies the temporal condition by construction. The result also shows that survival of death under the pattern theory is \emph{extrinsic}, in the sense of Corollary~3.4 of PID-02-a: whether $B$ is $x$ depends on whether there is a second realization, which is not a feature of the relation between $x$ and $B$. This extrinsicness will matter in Section~\ref{s:evid}.

\begin{corollary}\label{c:answers}
The gap problem, the duplication problem and the causal problem do not show survival of death to be incoherent under the pattern theory.
\end{corollary}

\begin{proof}
The gap problem is answered by Lemma~\ref{l:gap}. The duplication problem infers from the fact that two replicas would not be the person that one replica would not be either; the inference requires the only-$x$-and-$y$ principle, which Principle~\ref{p:persist} rejects for persons, as PID-02-a argues on independent grounds: the principle holds for kinds whose persistence relation cannot branch, and pattern continuity can branch. The causal problem assumes that the causal link must be immanent; Definition~\ref{d:cont} requires only that it be reliable, and Lemma~\ref{l:bridge} shows that the link of the model is reliable.
\end{proof}

Van Inwagen's manuscript deserves a separate word. A manuscript is an artefact, and on the conclusion of MET-07 an artefact persists through the continuity of its maintenance, not through the continuity of a pattern; the reassembled ship of Theseus is not the original ship. Van Inwagen is right that God could not restore Augustine's manuscript by creating a duplicate. The argument transfers to persons only if persons persist as artefacts do, and the pattern theory denies that they do.

\section{Models and further cases}\label{s:cases}

The model of Theorem~\ref{t:coh} is one among several. We apply the theory to the principal models of resurrection in the literature and to cases that test its limits.

\paragraph{Replica in another space.} Hick's resurrection world is a space with no spatial relation to ours, in which replicas of the dead appear \citep{Hick1976}. On the pattern theory the location of the replica is irrelevant; what matters is the causal link and the absence of a competitor. Hick did not specify a causal link, and on a theory that requires none, any exact replica would qualify, including one produced by accident. Definition~\ref{d:cont} requires that the replica depend reliably on the original, which God's knowledge can provide; Hick's model is coherent under the pattern theory once the dependence is added.

\paragraph{Simulacrum.} Van Inwagen's model, on which God preserves the body and leaves a simulacrum, is coherent under the pattern theory, and trivially so, since the realizer continues. It is also coherent under theories, such as animalism, that require the continuity of the organism. It is not needed under the pattern theory, and its main cost, which van Inwagen acknowledged, is that it makes the corpses we bury into substitutes.

\paragraph{Falling elevator.} On Zimmerman's model each simple at death has two immediate successors, one set making up the corpse and the other the body that is transported to the place of resurrection \citep{Zimmerman1999}. The corpse does not realize the person's pattern, so the model involves no branching of pattern continuity, and it is coherent under the pattern theory. It is designed to satisfy an immanent-causation requirement that the pattern theory does not impose.

\paragraph{Intermediate state.} Some traditions hold that the dead exist between death and resurrection. On the pattern theory a person exists in an intermediate state only if something realizes her pattern then. If the pattern can be realized by a non-physical system, a disembodied intermediate existence is coherent; PID-10 concludes that disembodied persons are possible in principle, and the pattern theory provides one way in which they could be. If, as MND-01 and MND-02 suggest, consciousness is the intrinsic nature of physical structure, the intermediate realizer must itself be physical, and an intermediate state is coherent only with a physical bearer. The theory does not require an intermediate state; Lemma~\ref{l:gap} makes it unnecessary.

\paragraph{Reincarnation.} A child who comes to realize, through a reliable process, a pattern connected by strong psychological connections to that of a dead person would, under the pattern theory, be that person, provided nothing else continues her. The reported cases do not come close to this: the reported memories are fragmentary and the child's values, skills and character are mostly her own. On the theory the question is one of degree of connection, and strong connection is lacking; whatever the cases show about the source of the child's information, they are not cases of the continuation of a pattern.

\paragraph{Death in dementia.} Suppose a person dies after years of dementia, and a resurrected body realizes the pattern she had before the illness. Is the result she? The later pattern is strongly connected to her earlier pattern, but only weakly to her pattern at death. Principle~\ref{p:persist} requires only that the later realization be linked to the earlier person by overlapping chains, and the resurrected realization is so linked to her earlier self, who is linked to her later self by the ordinary chains of a life. The resurrected person is therefore she, provided no other realization continues her. PID-08 concludes that survival does not require memory continuity, only the continuity of other psychological capacities, which supports the verdict for the intermediate period. What the case shows is that the restored pattern need not be the pattern at death.

\paragraph{Infants.} An infant who dies realizes a pattern too thin to satisfy Principle~\ref{p:persons}. If the pattern theory is the true theory of persons, an infant who dies before acquiring self-consciousness is not yet a person, and there is no person whose survival is in question; there may be the survival of an organism or of a pattern of another kind. Theologians who hold that infants are raised must either hold that the infant was a person on some other ground or that what is raised is the infant's developing life. The pattern theory forces the question but does not answer it.

\section{The pattern theory and the PID family}\label{s:pid}

The coherence result is relative to the pattern theory. We must say how that theory stands to the conclusions of the PID family, which concern what we in fact are.

\paragraph{Psychological continuity.} PID-01 concludes that personal identity over time consists in psychological continuity, with bodily continuity as a tie-breaker in cases of branching. The pattern theory agrees with the first clause and has no use for the second in the post-mortem case, since after death no body continues. If there were two post-mortem realizations, PID-01 and the pattern theory would agree that neither is the person, since neither is bodily continuous with her. PID-02-a's verdict on teletransportation follows from the same criterion, and Proposition~\ref{p:timing} shows that survival of death satisfies its temporal condition. The pattern theory is therefore an extension of the view of PID-01 and PID-02-a to cases with a temporal gap and a non-standard causal channel, not a rival to it.

\paragraph{Fission and what matters.} PID-03 and PID-06 conclude that in fission each successor has what matters in survival, and that psychological connectedness, not identity, carries most of what matters. If God produced two post-mortem realizations, the person would not survive, by Proposition~\ref{p:timing}, but each realization would have what matters. The distinction between survival and the possession by a successor of what matters is the distinction Section~\ref{s:evid} shows to be beyond evidence. That it is a distinction which, by PID-06, matters less than commonly assumed, is part of the reason the evidential result is less bleak than it might seem.

\paragraph{Animalism.} PID-05 concludes that you are a thinking organism and would not survive as a brain in a vat. If that is true, the pattern theory is not the true theory of us. The thing that thinks your thoughts is an organism; an organism persists through the continuity of its life, not of its psychology; and when the organism dies, you, on the animalist view, do not survive as a post-mortem realizer of your pattern \citep{Olson1997}. This is a genuine conflict, and we do not resolve it by weakening either conclusion. Three points define the relation. First, the coherence of survival under the pattern theory is not touched by the truth of animalism: a theory can be consistent and have models in which persons survive death even if we are not the persons that theory describes. Second, if PID-05 is true, the post-mortem realizer of the model of Theorem~\ref{t:coh} is not you, but it is a person, by Principle~\ref{p:persons}, and it has, by the lights of PID-06, much of what matters in your survival. A theist who accepts animalism and wants personal resurrection must adopt a model, such as van Inwagen's or Zimmerman's, on which the organism itself continues. Third, the choice between the pattern theory and animalism is not made here; it is made in the PID family. REL-08 establishes what follows for survival if the pattern theory is true and makes no claim that it is.

\paragraph{Uploading.} PID-14 concludes that uploading a mind creates a psychologically continuous successor, not the original person. If that conclusion rests on a requirement that the realizer be biological, or of the same kind as the original, it is a further reason to deny that the unrestricted pattern theory is true of us, since Definition~\ref{d:real} places no restriction on the substrate. The coherence result does not depend on the resolution: Theorem~\ref{t:coh} requires only that the post-mortem realizer be a body, and is compatible with requiring that it be a living organic body of the same kind as the original.

\section{The evidence}\label{s:evid}

We turn to the second question. We use the Bayesian account of confirmation defended in SCI-12. Let $S$ be the hypothesis that a given person $x$ survives death. Let $R$ be the hypothesis that at some time after $x$'s death some system realizes a pattern strongly connected to $x$'s.

\begin{definition}[Ante-mortem evidence]\label{d:ante}
Evidence is \emph{ante-mortem} for a subject if it consists of observations the subject makes, or receives through testimony, before her own death.
\end{definition}

\begin{definition}[Discriminating confirmation]\label{d:conf}
Evidence $E$ \emph{confirms} $H$ \emph{against} $H'$ if $P(E \mid H) > P(E \mid H')$. Evidence confirms $S$ \emph{as such} if it confirms $S$ against every hypothesis on which $x$'s pattern is realized after death and $x$ does not survive.
\end{definition}

Confirmation as such is what is at stake. Evidence that confirmed $R$ against the hypothesis that nothing realizes $x$'s pattern after death would confirm that something of $x$ continues, but not that $x$ does. Since $S$ entails $R$ but not conversely, the question whether evidence could confirm survival is the question whether it could favor $S$ over the hypotheses that agree with $S$ in making $R$ true and disagree in making $x$ fail to survive. Under the pattern theory there are two families of such hypotheses, corresponding to the two conditions that Principle~\ref{p:persist} adds to strong connection.

\begin{definition}[Matched rivals]\label{d:rivals}
Let $S$ be specified by a history $h$ of the world up to $x$'s death and a post-mortem realization $B$ that continues $x$. The \emph{duplicate rival} $D$ agrees with $S$ on $h$ and on $B$ and its states, and adds a second realization $B'$ pattern-continuous with $x$ that has no causal influence on the world of $h$ or on $B$. The \emph{deviant rival} $C$ agrees with $S$ on $h$ and on $B$'s states, but $B$ is produced by a process that does not depend reliably on $x$'s pattern.
\end{definition}

\begin{lemma}\label{l:fail}
Under the pattern theory, $x$ does not survive death under $D$ or under $C$.
\end{lemma}

\begin{proof}
Under $D$, $B$ and $B'$ are each pattern-continuous with $x$ and, since $B'$ has no causal influence on $B$ and conversely, neither is pattern-continuous with the other; so the continuity branches and by Principle~\ref{p:persist} neither is $x$. Under $C$, $B$ is not pattern-continuous with $x$, since the reliability clause of Definition~\ref{d:cont} fails, so $B$ is not $x$; and by construction nothing else continues $x$.
\end{proof}

\begin{assumption}[Locality of evidence]\label{a:local}
The probability of any ante-mortem observation, given a hypothesis about the world, depends only on what the hypothesis says about the world before the observer's death and about the states of things that can causally affect the observer before her death.
\end{assumption}

Assumption~\ref{a:local} says that what one can observe depends on what one can be affected by. It is the principle that makes observation evidence of anything at all, and it does not prejudge whether post-mortem realizations can affect the living; if they can, as mediums claim, Assumption~\ref{a:local} allows their states to bear on observations.

\begin{lemma}[Parity]\label{l:parity}
For every ante-mortem evidence $E$, $P(E \mid S) = P(E \mid D) = P(E \mid C)$.
\end{lemma}

\begin{proof}
By Definition~\ref{d:rivals}, $D$ agrees with $S$ on the history before death and on the states of $B$, and adds only $B'$, which has no causal influence on the world of the living or on $B$. $C$ agrees with $S$ on the history before death and on the states of $B$, and differs only in the process that produced $B$; since that process lies in the post-mortem world and its only effects there are the states of $B$, on which $C$ and $S$ agree, it makes no further difference to what can affect an ante-mortem observer. By Assumption~\ref{a:local}, the probability of $E$ is the same under all three.
\end{proof}

\begin{theorem}[Evidential inaccessibility]\label{t:inacc}
Under the pattern theory, no ante-mortem evidence confirms survival of death as such. For every ante-mortem $E$, the ratios $P(S \mid E) : P(D \mid E) : P(C \mid E)$ equal the prior ratios $P(S) : P(D) : P(C)$.
\end{theorem}

\begin{proof}
By Lemma~\ref{l:fail}, $D$ and $C$ are hypotheses on which $x$'s pattern is realized after death and $x$ does not survive. By Lemma~\ref{l:parity}, $P(E \mid S)$ is not greater than $P(E \mid D)$, so by Definition~\ref{d:conf} $E$ does not confirm $S$ against $D$, and hence does not confirm $S$ as such. By Bayes's theorem, $P(S \mid E)/P(D \mid E) = [P(E \mid S)/P(E \mid D)]\cdot[P(S)/P(D)] = P(S)/P(D)$, and likewise for $C$.
\end{proof}

\begin{corollary}[Ceiling]\label{c:ceiling}
For every ante-mortem $E$, $P(S \mid E) \leq P(S)/(P(S) + P(D) + P(C))$.
\end{corollary}

\begin{proof}
$S$, $D$ and $C$ are mutually exclusive, so $P(S \mid E) \leq P(S \mid E)/(P(S\mid E) + P(D \mid E) + P(C \mid E))$, and by Theorem~\ref{t:inacc} the right-hand side equals $P(S)/(P(S) + P(D) + P(C))$.
\end{proof}

Corollary~\ref{c:ceiling} states the result in its most useful form. Whatever evidence the living acquire, the probability of survival cannot exceed a bound fixed before any evidence is in, by the prior weight of survival against its matched rivals. Evidence may raise the probability that something continues the person's pattern after death, and in raising it may raise the probability of survival up to that bound; it cannot raise the bound. What decides between survival and its rivals is whether there is a second, causally isolated realization and whether the realization present was produced in the right way, and these are, by the extrinsicness of Proposition~\ref{p:timing}, facts about the post-mortem world that leave no trace in the ante-mortem one.

The result is specific to the pattern theory, and it is the theory's own structure that produces it. On a dualist theory, survival consists in the continued existence of a soul, and evidence of the soul's activity would, if obtainable, be evidence of survival; there is no matched rival of type $D$, since a soul cannot be duplicated. On an animalist theory, survival requires the continuation of the organism's life, which is in principle observable. On the pattern theory, by contrast, the facts that make the difference between survival and its matched rivals are, by design, facts about the relations among realizations, not about any realization's intrinsic states. That is why it can make survival coherent across a gap, and it is why no trace of survival can be distinguished from a trace of its absence.

\section{Mediumship, reincarnation and near-death experience}\label{s:evcases}

We apply Theorem~\ref{t:inacc} and Corollary~\ref{c:ceiling} to the three main bodies of putative evidence.

\paragraph{Mediumship.} The strongest mediumistic cases, such as those of Leonora Piper studied by James and Hodgson, involve the production by a medium, apparently in communication with a dead person, of information about that person's life that the medium had no normal means of acquiring \citep{Gauld1982}. Suppose, contrary to the skeptical reading, that the information came from a post-mortem realization of the dead person's pattern. Then the evidence confirms $R$ against the hypothesis that the information came from the living, by normal or by anomalous means. It does not confirm $S$ against $D$: the realization that communicates would communicate the same things whether or not a second realization existed elsewhere. Nor does it confirm $S$ against $C$: the realization would have the same memories, and would communicate them, whether or not it had been produced by a process that depended on the dead person's pattern. The discussion between \citet{Almeder1992} and \citet{Braude2003} concerns the comparison of $R$ with the hypothesis of information from the living, the so-called super-psi hypothesis. Theorem~\ref{t:inacc} shows that even its most favorable outcome for the survivalist would not settle survival. Nor would the result change if testimony about a medium's performance were plentiful and independent: REL-07 concludes that sufficiently many testimonies, independent given the truth of the report, can justify belief in a reported event of very low prior probability, and the reported event here, a medium's veridical communication, may be so established; what the testimony establishes is the event, not its interpretation as survival.

\paragraph{Reincarnation reports.} In the cases collected by \citet{Stevenson1974}, children report memories matching the life of a dead person. We argued in Section~\ref{s:cases} that the connections are too weak for pattern continuity. Even if they were strong, Theorem~\ref{t:inacc} would apply: whether the child is the dead person depends on whether some other realization continues the same pattern, and on the source of the child's states, and neither difference shows in the child's reports.

\paragraph{Near-death experience.} Reports of experiences during cardiac arrest \citep{vanLommel2001} concern subjects who did not die in the sense of Definition~\ref{d:death}, since their organisms resumed realizing their patterns. They bear, if on anything, on whether experience can occur while the brain's measurable activity is greatly reduced. That question is relevant to whether a person's pattern can be realized by something other than the brain, and so to $R$. It is not evidence of what happens after an irreversible death, and by Definition~\ref{d:ante} the subject's own experience, which she later reports, is ante-mortem evidence.

\paragraph{One's own case.} Before her death, a person has the evidence of her own experience, including any experiences she takes to be glimpses of what follows death. These are ante-mortem evidence, and Theorem~\ref{t:inacc} applies to them as to anything else. Whether she will be the one who wakes after death, rather than the first of two or a well-designed copy, is a matter on which nothing she can now observe bears.

\section{Objections and replies}

\begin{objection}[Persons as types]\label{o:types}
If a person is a pattern, then a person is an abstract type that can be multiply instantiated; persons so conceived cannot die, cannot be in two places at once without being two, and are not the concrete subjects we are.
\end{objection}

\begin{reply}
Principle~\ref{p:persons} identifies a person not with a pattern but with a concrete realization of one, and Principle~\ref{p:persist} identifies a person over time by a relation among realizations. A pattern can be realized twice; a person cannot be two realizations at once, which is why Principle~\ref{p:persist} contains the non-branching clause. Persons on this theory are as concrete as organisms, and die when their realizer ceases to realize their pattern; they differ from organisms in their persistence conditions, not in their category.
\end{reply}

\begin{objection}[Too many thinkers]\label{o:olson}
If a person is a realization distinct from the organism that realizes her, then the organism, which has the same brain, also thinks her thoughts. There are two thinkers wherever there is one person, and the person cannot know which she is \citep{Olson1997}.
\end{objection}

\begin{reply}
The objection is the principal argument for PID-05, and we do not claim to refute it; Section~\ref{s:pid} records that, if it succeeds, the pattern theory is not the true theory of us. But it is an objection to the truth of the theory, not to its coherence. A theory that identifies the person with the organism-as-realizing-the-pattern during life, and allows the person to be realized by another system afterwards, may be false, but no contradiction follows from it; the too-many-thinkers objection relies on the further premise that a thing's thinking cannot be derivative on another's, which the theory denies and which is not a logical truth. Theorem~\ref{t:coh} claims no more than coherence.
\end{reply}

\begin{objection}[Arbitrary rivals]\label{o:adhoc}
The duplicate and deviant rivals are skeptical hypotheses of the same kind as the hypothesis that the world was created five minutes ago. Any hypothesis faces such rivals, so Theorem~\ref{t:inacc} shows nothing special about survival.
\end{objection}

\begin{reply}
The rivals differ from global skeptical hypotheses in three respects. They are not constructed from the evidence: each is fixed in advance and matches $S$ on all ante-mortem evidence at once. They are not skeptical about the observable world: they agree with $S$ on everything before death and on everything any post-mortem realization does. And they are generated by the theory under which survival is coherent, not by a general doubt: the pattern theory makes survival depend on non-branching and on reliable causation, and the rivals are exactly the cases in which those conditions fail. A hypothesis about the past, by contrast, has no rival on which every trace is the same and the past is different, unless we invent one. The pattern theory supplies its own.
\end{reply}

\begin{objection}[Divine assurance]\label{o:reveal}
A truthful God who promised that the dead will be raised would not then produce duplicates or deviant copies. The rivals have a negligible prior given theism, and the ceiling of Corollary~\ref{c:ceiling} is close to one. Moreover, a revelation is itself evidence, available before death, that favors $S$ over its rivals.
\end{objection}

\begin{reply}
The first point concerns priors, and we accept it: Corollary~\ref{c:ceiling} allows the ceiling to be high for someone whose prior, on theological grounds, excludes the rivals. The confirmation of survival then comes from the theology, not from evidence about the dead. The second point would succeed only if a revelation available before death drew the distinction between being raised oneself and being continued by a single, rightly caused successor, since only such a revelation could have different likelihoods under $S$ and $D$. The promises of the scriptural traditions do not draw it; they are stated in terms neutral between the theories of persons. And the evidence that a given text or experience is a revelation at all is itself ante-mortem evidence, whose weight REL-21 assesses as low in the absence of cross-checking. We do not claim that a communication drawing the distinction is impossible, only that none is available; REL-03 gives a reason why a loving God would not give unambiguous revelation.
\end{reply}

\begin{objection}[Eschatological verification]\label{o:hick}
\citet{Hick1960} argued that statements about life after death are verifiable, since a person who survives will find them true. Survival is therefore confirmable after all.
\end{objection}

\begin{reply}
Hick's verification takes place after death, and our thesis concerns evidence available before it. Theorem~\ref{t:inacc} in fact extends further than we need: the post-mortem realization's own experiences are the same under $S$, $D$ and $C$, so even after death the survivor could not tell from her experience that she is not a duplicate or a deviant copy. Hick's argument shows that the post-mortem experience of a continuing life would confirm $R$; whether it would confirm survival depends on the theory of persons. REL-18 argues that religious statements are truth-apt though most of their truth-conditions cannot be verified; statements about survival are a central case.
\end{reply}

\section{What the result does not establish}\label{s:not}

The result does not establish that there is life after death, or that there is not. It does not establish that the pattern theory is true; Section~\ref{s:pid} records that PID-05 and, on one reading, PID-14 tell against it, and the coherence of survival under the theory is independent of its truth. It does not establish that evidence is irrelevant to the question. Evidence can bear on $R$, the hypothesis that one's pattern is realized after death, and by raising the probability of $R$ it can raise the probability of survival, up to the ceiling of Corollary~\ref{c:ceiling}. What evidence cannot do is move the ceiling, which is the sense in which it cannot confirm survival. The result does not claim that the existing mediumistic, reincarnation and near-death evidence supports $R$; that is a question for the assessment of the evidence, on which we take no position here, beyond noting that the near-death evidence is not evidence about the post-mortem period. Finally, the result is relative to the pattern theory. On a dualist or an animalist theory the matched rivals do not arise in the same form, and the evidential question would have to be asked again; the companion manuscript REL-08-b discusses the inaccessibility result in a form that does not depend on the details of the present framework.

\section{Conclusion}

Survival of death is coherent under a pattern theory of persons, on which a person is a concrete realization of a psychological pattern and persists through the reliably caused, non-branching continuation of that pattern. Under that theory a person may cease to be realized at death and be realized again later, through a causal link that runs through the knowledge of an agent, and such survival meets the temporal condition that governs teletransportation. The objections from the burned manuscript, from duplication and from the requirement of immanent causation presuppose persistence conditions the theory rejects. But the features that make survival coherent, its dependence on the absence of a second realization and on the right kind of cause, are invisible from before death. Evidence available to the living may bear on whether one's pattern will be realized again; it cannot favor one's surviving over one's being succeeded.

\begin{thebibliography}{99}\small
\bibitem[Almeder(1992)]{Almeder1992} Almeder, R. (1992). \emph{Death and Personal Survival: The Evidence for Life after Death}. Rowman \& Littlefield.
\bibitem[Aquinas(Cor)]{AquinasCor} Aquinas, T. \emph{Super I Epistolam ad Corinthios lectura}. Cited by chapter and lecture.
\bibitem[Baker(2000)]{Baker2000} Baker, L.~R. (2000). \emph{Persons and Bodies: A Constitution View}. Cambridge University Press.
\bibitem[Braude(2003)]{Braude2003} Braude, S.~E. (2003). \emph{Immortal Remains: The Evidence for Life after Death}. Rowman \& Littlefield.
\bibitem[Broad(1925)]{Broad1925} Broad, C.~D. (1925). \emph{The Mind and Its Place in Nature}. Kegan Paul.
\bibitem[Butler(1736)]{Butler1736} Butler, J. (1736). Of personal identity. Dissertation I appended to \emph{The Analogy of Religion}.
\bibitem[Corcoran(2006)]{Corcoran2006} Corcoran, K.~J. (2006). \emph{Rethinking Human Nature: A Christian Materialist Alternative to the Soul}. Baker Academic.
\bibitem[Gauld(1982)]{Gauld1982} Gauld, A. (1982). \emph{Mediumship and Survival: A Century of Investigations}. Heinemann.
\bibitem[Hick(1960)]{Hick1960} Hick, J. (1960). Theology and verification. \emph{Theology Today}, 17(1), 12--31.
\bibitem[Hick(1976)]{Hick1976} Hick, J. (1976). \emph{Death and Eternal Life}. Collins.
\bibitem[Hudson(2001)]{Hudson2001} Hudson, H. (2001). \emph{A Materialist Metaphysics of the Human Person}. Cornell University Press.
\bibitem[James(1898)]{James1898} James, W. (1898). \emph{Human Immortality: Two Supposed Objections to the Doctrine}. Houghton Mifflin.
\bibitem[Locke(1694)]{Locke1694} Locke, J. (1694). \emph{An Essay Concerning Human Understanding}, 2nd edn. Cited by book, chapter and section.
\bibitem[Olson(1997)]{Olson1997} Olson, E.~T. (1997). \emph{The Human Animal: Personal Identity Without Psychology}. Oxford University Press.
\bibitem[Parfit(1984)]{Parfit1984} Parfit, D. (1984). \emph{Reasons and Persons}. Clarendon Press.
\bibitem[Penelhum(1970)]{Penelhum1970} Penelhum, T. (1970). \emph{Survival and Disembodied Existence}. Routledge \& Kegan Paul.
\bibitem[Polkinghorne(2002)]{Polkinghorne2002} Polkinghorne, J. (2002). \emph{The God of Hope and the End of the World}. Yale University Press.
\bibitem[Shoemaker(1984)]{Shoemaker1984} Shoemaker, S. (1984). A materialist's account. In S. Shoemaker and R. Swinburne, \emph{Personal Identity}. Blackwell.
\bibitem[Steinhart(2014)]{Steinhart2014} Steinhart, E. (2014). \emph{Your Digital Afterlives: Computational Theories of Life after Death}. Palgrave Macmillan.
\bibitem[Stevenson(1974)]{Stevenson1974} Stevenson, I. (1974). \emph{Twenty Cases Suggestive of Reincarnation}, 2nd edn. University Press of Virginia.
\bibitem[Swinburne(1997)]{Swinburne1997} Swinburne, R. (1997). \emph{The Evolution of the Soul}, revised edn. Clarendon Press.
\bibitem[van Inwagen(1978)]{vanInwagen1978} van Inwagen, P. (1978). The possibility of resurrection. \emph{International Journal for Philosophy of Religion}, 9(2), 114--121.
\bibitem[van Lommel et~al.(2001)]{vanLommel2001} van Lommel, P., van Wees, R., Meyers, V. and Elfferich, I. (2001). Near-death experience in survivors of cardiac arrest: A prospective study in the Netherlands. \emph{The Lancet}, 358(9298), 2039--2045.
\bibitem[Wiggins(1980)]{Wiggins1980} Wiggins, D. (1980). \emph{Sameness and Substance}. Blackwell.
\bibitem[Zimmerman(1999)]{Zimmerman1999} Zimmerman, D.~W. (1999). The compatibility of materialism and survival: The \enquote{falling elevator} model. \emph{Faith and Philosophy}, 16(2), 194--212.
\end{thebibliography}

\end{document}
